
\SGARefTarget{sga7.sga7i.eviii.n2-3}{2.3}. \underline{Relation with ``Cartier duality''.}
The description given in \SGARefLink{sga7.sga7i.eviii.n2-2}{2.2} of the pairings defined by a
biextension, through the isomorphism theorem VII \SGARefLink{sga7.sga7i.evii.n3-6-5}{3.6.5} and the
homomorphism (\SGARefLink{sga7.sga7i.eviii.n2-2-1}{2.2.1}), is plainly not the definition formerly used for
abelian schemes, at a time when the very notion of biextension did not
yet exist.  That definition proceeds through Cartier duality, and we
now give its relation to the present one.

Consider again the homomorphism
\begin{equation}\tag{\SGARefTarget{sga7.sga7i.eviii.n2-3-1}{2.3.1}}
 \mathrm{Biext}^1(P,Q;G)
 \longrightarrow
 \operatorname{Hom}\bigl(P,\uExt^1(Q,G)\bigr)
\end{equation}
arising from the exact sequence (\SGARefLink{sga7.sga7i.eviii.n1-4}{1.4}).  It is in fact an isomorphism
when
\begin{equation}\tag{\SGARefTarget{sga7.sga7i.eviii.n2-3-2}{2.3.2}}
 \uHom(Q,G)=0.
\end{equation}
This condition holds in particular when $Q$ is an abelian scheme over
a scheme $S$ and $G$ is the multiplicative group $\underline G_m$
(VII \SGARefLink{sga7.sga7i.evii.n1-3-8}{1.3.8}).

There is also an evident ``restriction to $nQ$'' homomorphism
\[
 \uExt^1(Q,G)\longrightarrow\uExt^1(nQ,G),
\]
and hence a composite homomorphism
\begin{equation}\tag{\SGARefTarget{sga7.sga7i.eviii.n2-3-4}{2.3.4}}
\begin{split}
 \operatorname{Hom}\bigl(P,\uExt^1(Q,G)\bigr)
 &\longrightarrow
 \operatorname{Hom}\bigl(P,\uExt^1(nQ,G)\bigr)\\
 &\longrightarrow
 \operatorname{Hom}\bigl({}_nP,{}_n\uExt^1(nQ,G)\bigr),
\end{split}
\end{equation}
where the second arrow is restriction to $\,{}_nP$.  Composing with
(\SGARefLink{sga7.sga7i.eviii.n2-3-1}{2.3.1}), we obtain a homomorphism
\begin{equation}\tag{\SGARefTarget{sga7.sga7i.eviii.n2-3-5}{2.3.5}}
 \mathrm{Biext}^1(P,Q;G)
 \longrightarrow
 \operatorname{Hom}\bigl({}_nP,{}_n\uExt^1(nQ,G)\bigr).
\end{equation}

On the other hand, the exact sequence
\begin{equation}\tag{\SGARefTarget{sga7.sga7i.eviii.n2-3-6}{2.3.6}}
 0\longrightarrow{}_nQ\longrightarrow Q
 \xrightarrow{\ n\ }nQ\longrightarrow0
\end{equation}
gives rise to a long exact sequence of the $\uExt^i(-,G)$, and hence
to the short exact sequence
\begin{equation}\tag{\SGARefTarget{sga7.sga7i.eviii.n2-3-7}{2.3.7}}
 0\longrightarrow\uHom(Q,G)_n
 \longrightarrow\uHom({}_nQ,G)
 \xrightarrow{\ \partial\ }
 {}_n\uExt^1(nQ,G)
 \longrightarrow0,
\end{equation}
where, as usual, $E_n$ denotes $E/nE$.  We retain from this sequence
the canonical surjective homomorphism
\begin{equation}\tag{\SGARefTarget{sga7.sga7i.eviii.n2-3-8}{2.3.8}}
 \uHom({}_nQ,G)
 \xrightarrow{\hspace{3em}\partial\hspace{3em}}
 {}_n\uExt^1(nQ,G),
\end{equation}
which is an isomorphism if $\uHom(Q,G)_n=0$, and a fortiori if (\SGARefLink{sga7.sga7i.eviii.n2-3-2}{2.3.2})
holds.  From (\SGARefLink{sga7.sga7i.eviii.n2-3-8}{2.3.8}) we deduce a homomorphism
\begin{equation}\tag{\SGARefTarget{sga7.sga7i.eviii.n2-3-9}{2.3.9}}
\begin{split}
 \operatorname{Hom}\bigl({}_nP,\uHom({}_nQ,G)\bigr)
 &\simeq
 \operatorname{Hom}({}_nP\otimes{}_nQ,G)\\
 &\longrightarrow
 \operatorname{Hom}\bigl({}_nP,{}_n\uExt^1(nQ,G)\bigr),
\end{split}
\end{equation}
which is an isomorphism when (\SGARefLink{sga7.sga7i.eviii.n2-3-2}{2.3.2}) holds.

All the homomorphisms under consideration are functorial in $P,Q,G$.
The homomorphisms (\SGARefLink{sga7.sga7i.eviii.n2-3-5}{2.3.5}), (\SGARefLink{sga7.sga7i.eviii.n2-3-9}{2.3.9}), and (\SGARefLink{sga7.sga7i.eviii.n2-2-2}{2.2.2}) form the diagram
\begin{equation}\tag{\SGARefTarget{sga7.sga7i.eviii.n2-3-10}{2.3.10}}
\begin{tikzcd}[column sep=large,row sep=large]
 &\mathrm{Biext}^1(P,Q;G)
   \arrow[ddl,"\varphi^n"']
   \arrow[ddr,"(2.3.5)"]&\\
 &-&\\
 \operatorname{Hom}({}_nP\otimes{}_nQ,G)
   \arrow[rr,"(2.3.9)"]&&
 \operatorname{Hom}\bigl({}_nP,{}_n\uExt^1(nQ,G)\bigr),
\end{tikzcd}
\end{equation}
where the horizontal arrow (\SGARefLink{sga7.sga7i.eviii.n2-3-9}{2.3.9}) is an isomorphism if (\SGARefLink{sga7.sga7i.eviii.n2-3-2}{2.3.2})
holds.

\medskip
\noindent\underline{Proposition \SGARefTarget{sga7.sga7i.eviii.n2-3-11}{2.3.11}.}
\underline{The diagram (\SGARefLink{sga7.sga7i.eviii.n2-3-10}{2.3.10}) is anticommutative.}

When $\uHom(Q,G)=0$, this statement therefore recovers a
characterization of the homomorphism $\varphi^n$ of (\SGARefLink{sga7.sga7i.eviii.n2-2-2}{2.2.2}) in terms
of the homomorphisms (\SGARefLink{sga7.sga7i.eviii.n2-3-5}{2.3.5}) and (\SGARefLink{sga7.sga7i.eviii.n2-3-9}{2.3.9}).  This is the
characterization formerly used in the theory of abelian schemes,
where moreover $nQ=Q$.
